\documentclass[11pt]{article}

\usepackage[margin=1in]{geometry}
\usepackage[T1]{fontenc}
\usepackage{lmodern}
\usepackage{microtype}
\usepackage{amsmath,amssymb}
\usepackage[dvipsnames]{xcolor}
\usepackage{enumitem}
\usepackage[most]{tcolorbox}
\usepackage[hidelinks]{hyperref}

\hypersetup{
  pdftitle={Homework 1, MATH 327, Fall 2026},
  pdfauthor={Manish Patnaik},
  pdfsubject={MATH 327 homework problems}
}

\definecolor{courseblue}{HTML}{1F4E79}
\definecolor{lightblue}{HTML}{EAF2F8}
\definecolor{boxborder}{HTML}{B8CBDC}
\definecolor{darkgray}{HTML}{333333}

\setlength{\parindent}{0pt}
\setlength{\parskip}{0.55em}
\setlist[enumerate]{leftmargin=2em, itemsep=0.25em, topsep=0.3em}

\newtcolorbox{exercise}[1]{
  enhanced,
  colback=white,
  colframe=boxborder,
  colbacktitle=lightblue,
  coltitle=courseblue,
  fonttitle=\bfseries,
  title={Exercise #1},
  boxrule=0.6pt,
  arc=1mm,
  left=2.5mm,
  right=2.5mm,
  top=1.5mm,
  bottom=1.5mm,
  before skip=0.8em,
  after skip=0.8em
}

\newcommand{\chapterheading}[1]{%
  \par\vspace{0.5em}
  {\color{courseblue}\Large\bfseries #1}\par\vspace{0.35em}
  {\color{boxborder}\hrule height 0.6pt}\vspace{0.35em}
}

\newcommand{\sectionheading}[1]{%
  \vspace{0.25em}
  {\color{darkgray}\large\bfseries #1}\par
}

\begin{document}

\begin{center}
  \colorbox{lightblue}{%
    \begin{minipage}{0.92\textwidth}
      \vspace{0.6em}
      \begin{tabular*}{\textwidth}{@{\extracolsep{\fill}}lr}
        {\color{courseblue}\bfseries\LARGE Homework 1} &
        {\color{courseblue}\bfseries\LARGE MATH 327}\\[0.4em]
        \multicolumn{2}{l}{\color{darkgray}\bfseries Fall 2026}
      \end{tabular*}
      \vspace{0.5em}
    \end{minipage}}
\end{center}

\vspace{0.5em}

The following exercises are from Michael Artin's \emph{Algebra}, second
edition.

\chapterheading{Chapter 1: Matrices}

\sectionheading{Section 1: The Basic Operations}

\begin{exercise}{1.13}
A square matrix $A$ is \emph{nilpotent} if $A^k=0$ for some $k>0$.
Prove that if $A$ is nilpotent, then $I+A$ is invertible. Do this by
finding the inverse.
\end{exercise}

\sectionheading{Miscellaneous Problems}

\begin{exercise}{M.10}
Let $A,B$ be $m\times n$ and $n\times m$ matrices. Prove that
$I_m-AB$ is invertible if and only if $I_n-BA$ is invertible.

\textit{Hint.} Perhaps the only approach available to you at this time is
to find an explicit expression for one inverse in terms of the other. As a
heuristic tool, you could try substituting into the power series expansion
for $(1-x)^{-1}$. The substitution will make no sense unless some series
converge, and this needn't be the case. But any way to guess a formula is
permissible, provided that you check your guess afterward.
\end{exercise}

\chapterheading{Chapter 2: Groups}

\sectionheading{Section 3: Subgroups of the Additive Group of Integers}

\begin{exercise}{3.3}
\begin{enumerate}[label=\textbf{(\alph*)}]
  \item Define the greatest common divisor of a set
        $\{a_1,\ldots,a_n\}$ of $n$ integers. Prove that it exists, and
        that it is an integer combination of $a_1,\ldots,a_n$.
  \item Prove that if the greatest common divisor of
        $\{a_1,\ldots,a_n\}$ is $d$, then the greatest common divisor of
        $\{a_1/d,\ldots,a_n/d\}$ is $1$.
\end{enumerate}
\end{exercise}

\newpage
\sectionheading{Section 4: Cyclic Groups}

\begin{exercise}{4.3}
Let $a$ and $b$ be elements of a group $G$. Prove that $ab$ and $ba$ have
the same order.
\end{exercise}

\begin{exercise}{4.7}
Let $x$ and $y$ be elements of a group $G$. Assume that each of the
elements $x$, $y$, and $xy$ has order $2$. Prove that the set
$H=\{1,x,y,xy\}$ is a subgroup of $G$, and that it has order $4$.
\end{exercise}

\sectionheading{Section 5: Homomorphisms}

\begin{exercise}{5.3}
Let $U$ denote the group of invertible upper triangular $2\times2$
matrices $A=\begin{bmatrix}a&b\\0&d\end{bmatrix}$, and let
$\varphi\colon U\to\mathbb{R}^{\times}$ be the map that sends
$A\mapsto a^2$. Prove that $\varphi$ is a homomorphism, and determine its
kernel and image.
\end{exercise}

\begin{exercise}{5.6}
Determine the center of $\operatorname{GL}_n(\mathbb{R})$.

\textit{Hint.} You are asked to determine the invertible matrices $A$ that
commute with every invertible matrix $B$. Do not test with a general matrix
$B$. Test with elementary matrices.
\end{exercise}

\sectionheading{Section 6: Isomorphisms}

\begin{exercise}{6.7}
Let $H$ be a subgroup of $G$, and let $g$ be a fixed element of $G$. The
\emph{conjugate subgroup} $gHg^{-1}$ is defined to be the set of all
conjugates $ghg^{-1}$, with $h$ in $H$. Prove that $gHg^{-1}$ is a
subgroup of $G$.
\end{exercise}

\begin{exercise}{6.11}
Let $a$ be an element of a group $G$. Prove that if the set $\{1,a\}$ is
a normal subgroup of $G$, then $a$ is in the center of $G$.
\end{exercise}

\vfill
{\small\color{darkgray}Compiled \today.}

\end{document}
